% $Header: /home/vedranm/bitbucket/beamer/solutions/conference-talks/conference-ornate-20min.en.tex,v 90e850259b8b 2007/01/28 20:48:30 tantau $

\documentclass{beamer} %[compress]
% alternatively add handout or draft option

%\mode<handout>%{\setbeamercolor{background canvas}{bg=black!5}}
%\usepackage{pgfpages}
%\pgfpagesuselayout{4 on 1}[a4paper,border shrink=5mm,landscape]
% For handout printing, 4 slides on 1 page

\mode<presentation>
{
  \usetheme{Pittsburgh} % plain template Pittsburgh
  \useoutertheme{infolines} % bottom (author, title, place)
  \useoutertheme[subsection=false]{miniframes} % top (sections, frames)
  \usecolortheme{beaver} %blue color whale
	\setbeamertemplate{frametitle}[default][right] % frame titles on the right
  \setbeamercovered{transparent}
  % or whatever (possibly just delete it)
}

\setbeamersize{text margin left=1cm, text margin right=1cm}
%\setbeamertemplate{navigation symbols}{}
\usepackage[english]{babel}
\usepackage{booktabs}
\usepackage{comment}
% or whatever

\usepackage[cp1250]{inputenc}
% or whatever
\usepackage{subfigure} %subfigures

\usepackage{times}
%\usepackage{breqn}
\usepackage{amsmath}
\usepackage[T1]{fontenc}
\def\sym#1{\ifmmode^{#1}\else\(^{#1}\)\fi} %shortcut for Stata tables
% Or whatever. Note that the encoding and the font should match. If T1
% does not look nice, try deleting the line with the fontenc.

%\logo{\includegraphics[height=0.5cm]{logo.jpg}} % logo on every page



\title[Lecture 2: Examples] % (optional, use only with long paper titles)
{Lecture 2: Applied Meta-Analysis Examples}

\author[T.\ Havranek] % (optional, use only with lots of authors)
{ Tomas Havranek
}
% - Give the names in the same order as the appear in the paper.
% - Use the \inst{?} command only if the authors have different
%   affiliation.

\institute[Prague] % (optional, but mostly needed)
{
  	Charles University, CEPR, METRICS

  %\inst%
  %Economic Research Department\\
  %Czech National Bank
      %\inst{1}%
  %Research Department\\
  %Czech National Bank
  %\and
  %\inst{2}%
  %Institute of Economic Studies\\
  %Charles University, Prague
  
%\pgfdeclareimage[height=1cm]{logo}{logo.jpg} % logo only on title page
%\pgfuseimage{logo}
}% - Use the \inst command only if there are several affiliations.
% - Keep it simple, no one is interested in your street address.

\date[Meta in Econ and Finance] % (optional, should be abbreviation of conference name)
{Research Synthesis in Economics and Finance, \\ University of Canterbury\\[2ex]
  February 18, 2025
}%\\[4mm]% - Either use conference name or its abbreviation.
% - Not really informative to the audience, more for people (including
%   yourself) who are reading the slides online


%\subject{Intertemporal Substitution}

\begin{comment}
\AtBeginSection[]
{
  \begin{frame}<beamer>{Outline}
    \tableofcontents[currentsection]%,currentsubsection]
  \end{frame}
}
\end{comment}

\begin{document}

\begin{frame}
  \titlepage
\end{frame}


\begin{frame}{Example 1: Daylight saving time}
When using econometric analysis, researchers typically estimate the following model: 
\begin{equation*}\label{eq:dst}
ln\text{Consumption}_{t} = \alpha + \text{\emph{DST}}\cdot\text{Treatment effect}_{t} + \text{Controls}_{t} + \epsilon.
\end{equation*}
\vspace{-0.5cm}
\begin{itemize}
	\item \emph{Consumption}: the average energy consumption during time \emph{t} for a given hour, day, and year.
\item \emph{Treatment effect}: equals 1 for all hours when daylight saving time applies. 
\item \emph{Controls}: seasonality and holidays, weather, the intensity of sunlight, heterogeneity among consumption units, etc. 
\end{itemize}
\end{frame}


\begin{frame}{No consensus in the literature}
\begin{center}
		\includegraphics[width=0.87\textwidth]{_trend.pdf}
\end{center}
\end{frame}






\begin{frame}{Journals report smaller savings}
\begin{center}
		\includegraphics[width=0.87\textwidth]{_density.pdf}
\end{center}
\end{frame}

\begin{frame}{Publication bias? Probably not}
\begin{center}
		\includegraphics[width=0.87\textwidth]{_funnel.pdf}
\end{center}
\end{frame}



\begin{frame}{Funnel asymmetry tests}
If the ratio of the point estimate to its standard error has a $t$-distribution, then the two quantities should be independent:
\begin{equation*}
\mbox{estimate}_i = \underbrace{\beta}_{\mbox{true effect}} + \underbrace{\beta_0 SE_i}_{\mbox{publication bias}} + \mu_{i}.
\end{equation*}

%The equation is heteroscedastic. Weighted least squares yield 

%\begin{equation*} 
%t_i = \beta_0 + \beta \left(\frac{1}{SE_i}\right) + \vartheta_{i}.
%\end{equation*} 
The no.\ of obs.\ can be used as an instrument for SE.

\bigskip

\begin{tiny}
\begin{tabular*}{1\hsize}{@{\hskip\tabcolsep\extracolsep\fill}l*{7}{c}}
\toprule
                &\multicolumn{1}{c}{OLS} &\multicolumn{1}{c}{FE}  &\multicolumn{1}{c}{BE}&\multicolumn{1}{c}{Country}&\multicolumn{1}{c}{ME}  &\multicolumn{1}{c}{IV}  \\
\midrule
SE (\emph{publication bias})     &   -0.410 &   -1.217&   -0.410&   -0.496&   -0.449 &    0.226  \\
                &  (0.265)&  (0.790)&  (0.757)&  (0.805)  &  (0.688)  &  (1.088) \\
Constant (\emph{true effect}) &   -0.293\sym{\sym{***}}&   -0.222\sym{\sym{***}}&   -0.294\sym{\sym{***}}&   -0.278\sym{\sym{***}}&   -0.291\sym{\sym{***}}&   -0.445\sym{\sym{*}}  \\
                &(0.000778)& (0.0700)& (0.00812)& (0.0459)&(0.00731)&  (0.243) \\
\midrule
Observations    &      101               &      101               &      101               &      101&   101&    90\\
\bottomrule
\end{tabular*}
\end{tiny}
\end{frame}

\begin{frame}{Country heterogeneity}
\begin{center}
		\includegraphics[width=0.65\textwidth]{_countries.pdf}
\end{center}
\end{frame}

\begin{frame}{Method heterogeneity (just studies for the US)}
\begin{center}
		\includegraphics[width=0.65\textwidth]{_studies.pdf}
\end{center}
\end{frame}



\begin{frame}{What explains the differences in conclusions? (1)}
\begin{scriptsize}
\begin{tabular}{@{\hskip\tabcolsep\extracolsep\fill}p{0.2\hsize}p{0.7\hsize}}%p{0.05\hsize}p{0.05\hsize}}
\toprule
  \multicolumn{2}{l}{\textit{Data characteristics}} \\
    Data period & The number of years used in the estimation. \\
    Main estimate & = 1 if the estimate is preferred by the authors of the study. \\
		Hourly data & = 1 if the data are examined on hourly or higher than hourly granularity. \\
    Daily data & = 1 if the data are examined on a daily basis. \\
    Daylight hours & Average time between sunrise and sunset on the longest day for the country or region under examination. \\
    Europe & = 1 if European countries are examined.  \\
    USA   & = 1 if US data are examined.  \\
    		\midrule
    \multicolumn{2}{l}{\textit{Design of the analysis}} \\
    Regression analysis & = 1 if the primary study is based on regression analysis.  \\
    Simulation analysis &  = 1 if the study is based on simulation. \\
    Difference-in-diff. &  = 1 if the difference-in-differences approach is employed.  \\
    Residential cons. &  = 1 if only residential consumption is examined.  \\
    Lighting cons. & = 1 if total energy savings are reported as a result of lighting reduction.  \\
\bottomrule
\end{tabular}
\end{scriptsize}
\end{frame}

\begin{frame}{What explains the differences in conclusions? (2)}
\begin{scriptsize}
\begin{tabular}{@{\hskip\tabcolsep\extracolsep\fill}p{0.2\hsize}p{0.7\hsize}}%p{0.05\hsize}p{0.05\hsize}}
\toprule
    \multicolumn{2}{l}{\textit{Publication characteristics}} \\
		Publication year & The publication year of the study (base = 1970).  \\
    Journal article &  = 1 if the study was published in a peer-reviewed journal. \\
    Impact factor & The recursive RePEc impact factor of the outlet. \\
    Citations & The logarithm of the total number of citations of the study in Google Scholar. \\
\bottomrule
\end{tabular}
\end{scriptsize}
\end{frame}



\begin{frame}{Bayesian model averaging}
\begin{center}
		\includegraphics[width=1\textwidth]{bma_uip.pdf}
\end{center}
\end{frame}



  
\begin{frame}{Results}
  \begin{block}{Main Findings}
  \begin{enumerate}
  \item
The mean reported estimate suggests energy savings of 0.34\% due to DST.
\item
The results are driven by the method employed (simulation vs.\ regression).
    \item
Energy savings are larger for countries farther away from the equator.
  \end{enumerate}
  \end{block}

\vskip0pt plus.5fill
    \begin{block}{Project Website}
      \href{http://www.meta-analysis.cz/dst}{www.meta-analysis.cz/dst}
    \end{block}
\end{frame}

% All of the following is optional and typically not needed. 
%\appendix
%\section<presentation>*{\appendixname}
%\subsection<presentation>*{For Further Reading}

\begin{frame}%[allowframebreaks]
  \frametitle<presentation>{For further reading}
    
  \begin{thebibliography}{10}
    
  \beamertemplatebookbibitems
  % Start with overview books.
    
    \bibitem[{Hunter \& Schmidt(1990)}]{Hunter1990}
Stanley, T.\ D.\ \& C. Doucouliagos (2012): {{Meta-Regression Analysis in Economics and Business}.}\newblock Routledge, 1st. edition.
  
  \beamertemplatearticlebibitems
  % Followed by interesting articles. Keep the list short. 
    
   \bibitem[{Stanley(2008)}]{Stanley2010}
Brodeur, A., M. Sangnier \& Y. Zylberg (2016): \\ {{Star Wars: The Empirics Strike Back}.} \newblock \emph{AEJ Applied} 8(1): pp. 1--32. 
    
    \bibitem[{Stanley(2001)}]{Stanley2001}
Havranek, T. (2015): {{Measuring Intertemporal Substitution: The Importance of Method Choices and Publication Bias}.}
\newblock \emph{JEEA} 13(6): pp.\ 1180--1204.
    \end{thebibliography}
\bigskip
Reading list on RePEc: Google ``meta-analysis in economics.'' %\href{http://ideas.repec.org/k/metaana.html}{ideas.repec.org/k/metaana.html}
    \end{frame}










\begin{frame}{Example 2: Interest rates and house prices}
\begin{figure}
	\centering
		\includegraphics[width=0.9\textwidth]{mean_irf_include}
\end{figure}
\end{frame}




\begin{frame}{Converting graphs to data}
We measure pixel coordinates using WebPlotDigitizer:
\smallskip
\smallskip
\begin{figure}
	\centering
		\includegraphics[width=1.0\textwidth]{measure.png}
\end{figure}

\end{frame}




\begin{frame}{Literature search}
\begin{figure}
	\centering
		\includegraphics[width=0.7\textwidth]{prisma.png}
\end{figure}
\end{frame}



\begin{frame}{Histograms for different horizons}
\begin{figure}
	\centering
		\includegraphics[width=0.9\textwidth]{hist}
\end{figure}
\end{frame}








\begin{frame}{Funnel asymmetry test (Card \& Krueger)}
In the absence of publication bias the funnel should be symmetrical:
\begin{equation*}
\mbox{estimate}_i = \underbrace{\beta}_{\mbox{true effect}} + \underbrace{\beta_0 SE_i}_{\mbox{publication bias}} + \mu_{i}.
\end{equation*}

The equation is heteroscedastic. Weighted least squares yield 

\begin{equation*} 
t_i = \beta_0 + \beta \left(\frac{1}{SE_i}\right) + \vartheta_{i}.
\end{equation*} 
The no.\ of obs.\ can be used as an instrument for SE.
\end{frame}






\begin{frame}{Publication bias for longer horizons?}
\begin{figure}
	\centering
		\includegraphics[width=0.9\textwidth]{funnels}
\end{figure}
\end{frame}


\begin{frame}{Corroborated by funnel asymmetry tests}
Note that the slope is negative, the intercept is smaller than the reported mean effect:

\smallskip
\smallskip

\begin{figure}
	\centering
		\includegraphics[width=1.05\textwidth]{fat.png}
\end{figure}
\end{frame}

\begin{frame}{Nonliner tests also imply bias}
Note that the corrected effect is smaller than the reported mean effect:

\smallskip
\smallskip
\begin{figure}
	\centering
		\includegraphics[width=1.05\textwidth]{fat2.png}
\end{figure}
\end{frame}




\begin{frame}{Mean impulse response beyond publication bias}
\begin{figure}
	\centering
		\includegraphics[width=0.9\textwidth]{beyond_bias}
\end{figure}
\end{frame}




\begin{frame}{Cross-country heterogeneity}
\begin{figure}
	\centering
		\includegraphics[width=0.9\textwidth]{country_response}
\end{figure}
\end{frame}


\begin{frame}{Estimation context}
\begin{itemize}
	\item \textbf{Data characteristics}: monthly, panel, length, age
		\item \textbf{Specification characteristics}: foreign IR, credit, consumption, res.\ invest., money supply, exch.\ rate, long-run IR, lags
		\item \textbf{Estimation characteristics}: BVAR, sign restr.\ HP, sign restr.\ other, nonrecursive
		\item \textbf{Publication characteristics}: citations, impact, published
				\item \textbf{Country characteristics}: mortgage-to-GDP, floating, maturity, IR, spread, prolonged high HP, PTI, crisis, prolonged low IR, tourism, income, inflation, popul.\ growth, permits, ownership
		
\end{itemize}

\end{frame}



\begin{frame}{Bayesian model averaging}
\begin{figure}
	\centering
		\includegraphics[width=0.85\textwidth]{bma}
\end{figure}
\end{frame}




\begin{frame}{Best practice}
We construct a hypothetical ideal study: no publication bias \\ (SE = 0), recent data, long series, nonrecursive identification, published, good journal, highly cited.

\bigskip

\begin{figure}
	\centering
		\includegraphics[width=1\textwidth]{best.png}
\end{figure}
\end{frame}




\begin{frame}{Impulse response implied by best practice}
\begin{figure}
	\centering
		\includegraphics[width=0.9\textwidth]{best_practice_v2}
\end{figure}
\end{frame}






\begin{frame}{Results}
  \begin{block}{Main Findings}
  \begin{enumerate}
  \item
Among 237 impulse responses from 37 studies, the mean decrease in house prices is 1.2\% two years after a 1 pp.\ increase in the policy rate.
\item
The mean is exaggerated because of publication bias.
    \item
Stronger transmission for i) more developed markets,\\ ii) flatter yield curve, iii) house price boom.
  \end{enumerate}
  \end{block}

\vskip0pt plus.5fill
    \begin{block}{Project Website}
      \href{http://www.meta-analysis.cz/house_prices}{www.meta-analysis.cz/house\_prices}
    \end{block}
\end{frame}




\begin{comment}
% All of the following is optional and typically not needed. 
%\appendix
%\section<presentation>*{\appendixname}
%\subsection<presentation>*{For Further Reading}

\begin{frame}%[allowframebreaks]
  \frametitle<presentation>{Recent meta papers of our team}
    
  \begin{thebibliography}{10}
  %\beamertemplatebookbibitems
  % Start with overview books.
   % \bibitem[{Hunter \& Schmidt(1990)}]{Hunter1990}
%Stanley, T.\ D.\ \& C. Doucouliagos (2012): {{Meta-Regression Analysis in Economics and Business}.}\newblock Routledge, 1st. edition.
  \beamertemplatearticlebibitems
  % Followed by interesting articles. Keep the list short. 
 
	    \bibitem[{Hunter \& Schmidt(1990)}]{Hunter1990}
Irsova, Z., P.\ Bom, T.\ Havranek, \& H.\ Rachinger (2024): {{Spurious Precision of Meta-Analysis of Observational Research}.}\newblock \emph{Nature Human Behaviour}, revisions requested.
	
	   \bibitem[{Stanley(2008)}]{Stanley2010}
Havranek, T., Z.\ Irsova, L.\ Laslopova, \& O.\ Zeynalova (2024): {{Publication and Attenuation Biases in Measuring Skill Substitution}.}\newblock \emph{Review of Economics and Statistics}, forthcoming.

    \bibitem[{Stanley(2001)}]{Stanley2001}
Opatrny, M.\, T.\ Havranek, Z.\ Irsova, \& M.\ Scasny (2024): {{Publication Bias and Model Uncertainty in Measuring the Effect of Class Size on Achievement}.}
\newblock \emph{Journal of Labor Economics}, revisions requested.
    


\end{comment}



\begin{frame}{Example 3: Working while in school}
\begin{itemize}
	\item  Hard to influence commonly cited determinants of educational outcomes: ability, ethnicity, gender, parental
education and affluence
\item But most students can decide whether or not to work
\item How does that influence educational outcomes?
\item Should you motivate your kids to work while in school?
\item Mixed empirical results: from large negative to positive estimates	
	\item 861 estimates from 69 studies
\end{itemize}
\end{frame}


\begin{frame}{No consensus in the literature}
\begin{figure}
	\centering
		\includegraphics[width=0.9\textwidth]{trend.pdf}
	\label{fig:trend}
\end{figure}

\end{frame}




\begin{frame}{Negative for all countries except Germany}
\begin{figure}
	\centering
		\includegraphics[width=1.00\textwidth]{countries.pdf}
	\label{fig:countries}
\end{figure}

\end{frame}



\begin{frame}{Work intensity clearly matters}
\begin{figure}
	\centering
		\includegraphics[width=0.9\textwidth]{pattern3.pdf}
	\label{fig:pattern3}
\end{figure}

\end{frame}




\begin{frame}{Endogeneity control is important}
\begin{figure}
	\centering
		\includegraphics[width=0.9\textwidth]{pattern5.pdf}
	\label{fig:pattern5}
\end{figure}

\end{frame}




\begin{frame}{Little publication bias on average}
\begin{figure}
	\centering
		\includegraphics[width=0.9\textwidth]{maer2022/funnel.pdf}
	\label{fig:funnel}
\end{figure}

\end{frame}




\begin{frame}{No bias in studies addressing endogeneity}
\begin{figure}
	\centering
		\includegraphics[width=1.00\textwidth]{maer2022/iv.pdf}
	\label{fig:iv}
\end{figure}

\end{frame}





\begin{frame}{Strong bias in studies ignoring endogeneity}
\begin{figure}
	\centering
		\includegraphics[width=0.97\textwidth]{maer2022/ols.pdf}
	\label{fig:ols}
\end{figure}

\end{frame}







\begin{frame}{Model uncertainty; endogeneity control is key}
\begin{figure}
	\centering
		\includegraphics[width=0.82\textwidth]{bmax.pdf}
	\label{fig:bma}
\end{figure}

\end{frame}




\begin{frame}{Closer look: work intensity, measurement}
\begin{figure}
	\centering
		\includegraphics[width=1\textwidth]{densities.pdf}
	\label{fig:densities}
\end{figure}

\end{frame}



\begin{frame}{Closer look: endogeneity, panel data, Germany}
\begin{figure}
	\centering
		\includegraphics[width=1.00\textwidth]{densities1.pdf}
	\label{fig:densities1}
\end{figure}

\end{frame}





\begin{frame}{Best practice estimates}
\begin{figure}
	\centering
		\includegraphics[width=1.00\textwidth]{best.pdf}
	\label{fig:best}
\end{figure}

\end{frame}




  
\begin{frame}{Results}
  \begin{block}{Main Findings}
  \begin{enumerate}
  \item
Publication and endogeneity biases interact.
\item
The corrected effect is negative but economically zero.
    \item
Effects negative for high-intensity work, positive for Germany.
  \end{enumerate}
  \end{block}

\vskip0pt plus.5fill
    \begin{block}{Project Website}
      \href{http://www.meta-analysis.cz/students}{www.meta-analysis.cz/students}
    \end{block}
\end{frame}





\end{document}


